\chapter*{Závěr}
\addcontentsline{toc}{chapter}{Závěr}

% Shrnutí cílů a jejich naplnění, přínosy redesignu pro organizaci i budoucí zákazníky, reflexe postupu a metod, možnosti pokračování projektu

V rámci diplomové práce se podařilo úspěšně navrhnout a na uživatelích ověřit nový uživatelský design aplikace pro správu pracovního prostředí. S využitím metod uživatelsky orientovaného designu byl vytvořen komplexní návrh, který je nyní připraven k technické implementaci.

Během řešení práce byly zodpovězeny všechny stanovené výzkumné otázky:

\begin{itemize}
    \item \textbf{K otázce hlavních bariér v použitelnosti (VO1)} analýza a výzkum ukázaly, že uživatelé nejvíce naráželi na nekoherentní navigaci, absenci logické vizuální hierarchie a roztříštěnost ovládacích prvků napříč různými moduly, které vznikaly organicky a bez jednotného zadání.
    \item \textbf{Pro zajištění vizuální konzistence a budoucího škálování (VO2)} byl vytvořen ucelený design systém. Tento systém definuje nejen základní vizuální styl, ale funguje jako standardizovaná knihovna komponent, která umožní vývojovému týmu efektivně přidávat nové funkcionality bez narušení jednotného vzhledu a chování aplikace.
    \item \textbf{Zhodnocení vlivu redesignu na kognitivní zátěž (VO3)} proběhlo prostřednictvím uživatelského testování interaktivních prototypů. Výsledky testování použitelnosti jasně prokázaly, že nové rozhraní výrazně zvyšuje intuitivnost používání – uživatelé dokončovali zadané úkoly rychleji, s menší chybovostí a subjektivně vnímali práci se systémem jako méně zatěžující.
\end{itemize}

Výstupem diplomové práce je soubor vzájemně provázaných artefaktů, které naplňují definované cíle: výzkumná zpráva shrnující úvodní analýzu, ucelený design systém, interaktivní prototypy redesignovaných klíčových obrazovek připravené k implementaci, report o testování použitelnosti a grafický manuál definující metodický rámec pro další rozvoj.

Přínosem práce je vytvoření robustního a testováním ověřeného základu, který transformuje funkčně orientovaný interní nástroj do podoby moderního, přehledného a škálovatelného produktu. Z hlediska budoucího rozvoje se doporučuje bezproblémové předání (hand-off) vytvořených artefaktů vývojovému týmu k implementaci. Následný výzkum by se měl zaměřit na sběr kvantitativních dat z reálného provozu a případné rozšíření design systému o pokročilé standardy digitální přístupnosti.
